% Options for packages loaded elsewhere
\PassOptionsToPackage{unicode}{hyperref}
\PassOptionsToPackage{hyphens}{url}
\documentclass[
]{article}
\usepackage{xcolor}
\usepackage[margin=1in]{geometry}
\usepackage{amsmath,amssymb}
\setcounter{secnumdepth}{-\maxdimen} % remove section numbering
\usepackage{iftex}
\ifPDFTeX
  \usepackage[T1]{fontenc}
  \usepackage[utf8]{inputenc}
  \usepackage{textcomp} % provide euro and other symbols
\else % if luatex or xetex
  \usepackage{unicode-math} % this also loads fontspec
  \defaultfontfeatures{Scale=MatchLowercase}
  \defaultfontfeatures[\rmfamily]{Ligatures=TeX,Scale=1}
\fi
\usepackage{lmodern}
\ifPDFTeX\else
  % xetex/luatex font selection
\fi
% Use upquote if available, for straight quotes in verbatim environments
\IfFileExists{upquote.sty}{\usepackage{upquote}}{}
\IfFileExists{microtype.sty}{% use microtype if available
  \usepackage[]{microtype}
  \UseMicrotypeSet[protrusion]{basicmath} % disable protrusion for tt fonts
}{}
\makeatletter
\@ifundefined{KOMAClassName}{% if non-KOMA class
  \IfFileExists{parskip.sty}{%
    \usepackage{parskip}
  }{% else
    \setlength{\parindent}{0pt}
    \setlength{\parskip}{6pt plus 2pt minus 1pt}}
}{% if KOMA class
  \KOMAoptions{parskip=half}}
\makeatother
\usepackage{longtable,booktabs,array}
\newcounter{none} % for unnumbered tables
\usepackage{calc} % for calculating minipage widths
% Correct order of tables after \paragraph or \subparagraph
\usepackage{etoolbox}
\makeatletter
\patchcmd\longtable{\par}{\if@noskipsec\mbox{}\fi\par}{}{}
\makeatother
% Allow footnotes in longtable head/foot
\IfFileExists{footnotehyper.sty}{\usepackage{footnotehyper}}{\usepackage{footnote}}
\makesavenoteenv{longtable}
\usepackage{graphicx}
\makeatletter
\newsavebox\pandoc@box
\newcommand*\pandocbounded[1]{% scales image to fit in text height/width
  \sbox\pandoc@box{#1}%
  \Gscale@div\@tempa{\textheight}{\dimexpr\ht\pandoc@box+\dp\pandoc@box\relax}%
  \Gscale@div\@tempb{\linewidth}{\wd\pandoc@box}%
  \ifdim\@tempb\p@<\@tempa\p@\let\@tempa\@tempb\fi% select the smaller of both
  \ifdim\@tempa\p@<\p@\scalebox{\@tempa}{\usebox\pandoc@box}%
  \else\usebox{\pandoc@box}%
  \fi%
}
% Set default figure placement to htbp
\def\fps@figure{htbp}
\makeatother
\setlength{\emergencystretch}{3em} % prevent overfull lines
\providecommand{\tightlist}{%
  \setlength{\itemsep}{0pt}\setlength{\parskip}{0pt}}
\usepackage{bookmark}
\IfFileExists{xurl.sty}{\usepackage{xurl}}{} % add URL line breaks if available
\urlstyle{same}
\hypersetup{
  pdftitle={Understanding the Primary Socioeconomic Drivers of State Test Scores{,} Grades 3-8},
  pdfauthor={Mohammed Abdalla, Miriam Luka, Jayant Kumar, Kashif Awan},
  hidelinks,
  pdfcreator={LaTeX via pandoc}}

\title{Understanding the Primary Socioeconomic Drivers of State Test
Scores, Grades 3-8}
\author{Mohammed Abdalla, Miriam Luka, Jayant Kumar, Kashif Awan}
\date{}

\begin{document}
\maketitle

\section{Introduction}\label{introduction}

Standardized test scores are one of the statistics available for nearly
every public school district in the country, making them a common
measure to assess whether our education system is serving students well.
Recent reporting on a generation-long decline in reading scores, paired
with proposed cuts to federal Department of Education funding, has drawn
increased attention to the matter of where funding should be allocated
to improve test scores and more generally educational outcomes. Before
discussing funding allocation, policymakers require a clear
understanding of the data itself and the potential shared
characteristics among school districts with lower test scores.

Socioeconomic status is a characteristic consistently mentioned in this
conversation. Built off the correlations indicating that districts with
higher incidents of poverty, and less educational attainment among
adults tend to have lower-performing school test scores. But
socioeconomic status bundles together several distinct conditions:
income, education, employment, family structure, and use of public
assistance (such as SNAP benefits utilization, etc.). Understanding and
quantifying the relationship between socioeconomic status indicators and
district state test scores with increased specificity is crucial when
attempting to improve outcomes. A policymaker deciding how to allocate
funding should understand not just that socioeconomic status is
important, but which dimensions of it are most strongly associated with
test performance, and how large those associations are in practice.

We approach this question as data consultants for a coalition of school
boards seeking to better understand the primary drivers of achievement
gaps in their districts. Our research question is purely descriptive:
what is the relationship between district-level socioeconomic
characteristics and average student test performance across U.S. school
districts? We do not claim that any of these characteristics cause test
scores to rise or fall, only that describing the strength and shape of
these relationships is a necessary first step before anyone can design a
targeted, evidence-based funding policy or a causal study of one.

\section{Data}\label{data}

\subsection{Source and Provenance}\label{source-and-provenance}

We use the Stanford Education Data Archive (SEDA), version 6.0 (Reardon
et al., Stanford University), a publicly available dataset built from
over 480 million standardized test records submitted by U.S. states
under federal reporting requirements. SEDA converts each state's own
assessments onto a common scale by linking them to the National
Assessment of Educational Progress (NAEP), producing test scores that
are comparable across states and years. We combine two SEDA files, both
at the school-district level:

\begin{itemize}
\item
  \textbf{\texttt{seda\_geodist\_annual\_cs\_6.0}} (achievement):
  average math and reading/language-arts performance for grades 3--8,
\item
  \textbf{\texttt{seda\_cov\_geodist\_annual\_6.0}} (covariates):
  district-level socioeconomic and demographic characteristics compiled
  from the American Community Survey and other federal sources.
\end{itemize}

For additional information regarding the methodology the SEDA utilized
in collecting this data, as well as the files themselves, please visit
their website at:
\url{https://edopportunity.org/opportunity/data/downloads/\#documentation-6}

\subsection{Data Wrangling}\label{data-wrangling}

From the achievement file, we keep only the \texttt{all\ students}
aggregate (filtering on col \texttt{subcat} = \texttt{all}), which
reports the achievement data in aggregate instead of by subgroups such
as race or gender. Furthermore for both our achievement file and
covariate file we will filter for \texttt{year} = \textbf{2018}, as this
is the most recent pre-pandemic year available to completely avoid any
issues that may have resulted from COVID. Resulting in each row
representing one achievement observation per district for 2018. We then
join our data to the 2018 covariate file using the district identifier
\texttt{sedalea}, which is common to both files. Districts that appear
in only one file are dropped via an inner join, and rows missing any of
our six predictor variables are also dropped.

Before dropping missing values, the 2018 achievement file has 11,139
district-level \textbf{all students} observations and the 2018 covariate
file has 12,690. The inner join drops \textbf{1,551 covariate districts
(12.2\%)} that have no 2018 achievement estimate in any subgroup, which
is SEDA's suppression rule for districts with too few tested students.
The loss runs one way only: every district with an achievement estimate
has a matching covariate record, so the join itself discards nothing
from the achievement side. A further \textbf{95 matched districts} are
missing one or more of our six socioeconomic predictors and are dropped
as well, leaving a population dataset of \textbf{11,044} observations,
one per school district. (For reference, 324 districts in the covariate
file carry no socioeconomic data at all; 95 of those are among the
11,139 that matched, and are the ones removed here.) As discussed in our
previous proposal, districts are not drawn independently, districts
neighboring each other likely have similar labor markets, demographics,
and state funding resources which would naturally result in iid being
violated. We will address this directly with a random sampling method
rather than fitting on the full population with a resulting sample size
of 1,000.

\section{Data Prep}\label{data-prep}

\section{Sampling Methodology}\label{sampling-methodology}

From the 2018 SEDA, whose all-subgroups population file gives one record
per U.S. public school district; we restrict the frame to districts with
positive enrollment and complete data on the outcome (average test score
composite), the six socioeconomic predictors, total enrollment, and the
race and ethnicity shares. From that cleaned frame we draw a simple
random sample of 1,000 districts without replacement under a fixed seed,
which satisfies the IID assumption directly since every district is
selected independently and with equal probability. We then split the
sample before any analysis --- 300 districts (30\%) for exploration and
700 (70\%) held for confirmation --- so that variable selection and
diagnostics never touch the data used for final estimation, and five
automated checks verify sample size, missingness, positive enrollment,
unique identifiers, and that the two splits sum to 1,000. We used IID
because we wanted to use the more interesting models from class.

\section{Operationalization and data
visualisation}\label{operationalization-and-data-visualisation}

\section{Operationalization}\label{operationalization}

Two concepts have to become variables: \emph{student achievement} and
\emph{socioeconomic status}. Neither mapping is obvious, and both admit
alternatives we rejected for stated reasons.

{\def\LTcaptype{none} % do not increment counter
\begin{longtable}[]{@{}
  >{\raggedright\arraybackslash}p{(\linewidth - 6\tabcolsep) * \real{0.2500}}
  >{\raggedright\arraybackslash}p{(\linewidth - 6\tabcolsep) * \real{0.2500}}
  >{\raggedright\arraybackslash}p{(\linewidth - 6\tabcolsep) * \real{0.2500}}
  >{\raggedright\arraybackslash}p{(\linewidth - 6\tabcolsep) * \real{0.2500}}@{}}
\toprule\noalign{}
\begin{minipage}[b]{\linewidth}\raggedright
\end{minipage} & \begin{minipage}[b]{\linewidth}\raggedright
Variable
\end{minipage} & \begin{minipage}[b]{\linewidth}\raggedright
Column
\end{minipage} & \begin{minipage}[b]{\linewidth}\raggedright
Definition
\end{minipage} \\
\midrule\noalign{}
\endhead
\bottomrule\noalign{}
\endlastfoot
\textbf{Outcome (Y)} & Test performance & \texttt{cs\_mn\_avg\_ol} &
District mean score, math and reading, grades 3--8, in national SD
units \\
\textbf{Predictors (X)} & Income & \texttt{lninc50all} & Log of median
household income \\
& Educational attainment & \texttt{baplusall} & \% of adults with a
bachelor's degree or higher \\
& Poverty & \texttt{povertyall} & \% of residents below the poverty
line \\
& Unemployment & \texttt{unempall} & \% unemployment rate \\
& SNAP receipt & \texttt{snapall} & \% of households receiving SNAP
benefits \\
& Single motherhood & \texttt{single\_momall} & \% of households headed
by a single mother \\
\end{longtable}
}

\subsection{The outcome: district test
performance}\label{the-outcome-district-test-performance}

SEDA offers two candidate outcome columns. We use
\texttt{cs\_mn\_avg\_ol} (the ordinary least-squares estimate) rather
than \texttt{cs\_mn\_avg\_eb} (an empirical-Bayes estimate that shrinks
small districts toward the state mean). Shrinkage would improve
precision but at the cost of \emph{deliberately} pulling small districts
toward a common value exactly the between-district variation we are
trying to describe. For a descriptive study the unshrunken estimate is
the honest choice.

{\def\LTcaptype{none} % do not increment counter
\begin{longtable}[]{@{}
  >{\raggedright\arraybackslash}p{(\linewidth - 4\tabcolsep) * \real{0.5139}}
  >{\raggedleft\arraybackslash}p{(\linewidth - 4\tabcolsep) * \real{0.0972}}
  >{\raggedright\arraybackslash}p{(\linewidth - 4\tabcolsep) * \real{0.3889}}@{}}
\toprule\noalign{}
\begin{minipage}[b]{\linewidth}\raggedright
Statistic
\end{minipage} & \begin{minipage}[b]{\linewidth}\raggedleft
Value
\end{minipage} & \begin{minipage}[b]{\linewidth}\raggedright
Reference
\end{minipage} \\
\midrule\noalign{}
\endhead
\bottomrule\noalign{}
\endlastfoot
Districts & 300 & -- \\
Distinct values & 300 & 300 if no ties \\
Districts within 0.05 of the minimum & 1 & high value = floor effect \\
Districts within 0.05 of the maximum & 3 & high value = ceiling
effect \\
Mean & -0.005 & -- \\
SD & 0.380 & -- \\
Min & -1.112 & -- \\
25th pct & -0.265 & -- \\
Median & -0.027 & -- \\
75th pct & 0.230 & -- \\
Max & 1.134 & -- \\
Skewness & 0.193 & 0 if symmetric \\
Kurtosis & 3.210 & 3 if normal \\
\end{longtable}
}

The table shows the variable needs no transformation. All 300 districts
have a distinct score and none pile up at either extreme, so nothing is
censored. Skewness 0.19 and kurtosis 3.21 sit near their symmetric and
normal reference values. Symmetry is not an OLS requirement the
large-sample model assumes nothing about the marginal distribution of
\(Y\) but it does mean nothing suggests the outcome behaves
multiplicatively, so \textbf{we leave \(Y\) untransformed}. The scale is
already interpretable: a coefficient reads directly as a change in
standard deviations.

One limitation to state plainly. \texttt{cs\_mn\_avg\_ol} is an
\emph{estimate}, and SEDA publishes a standard error for each district.
About 1.9\% of the outcome's variance is measurement error rather than
genuine difference between districts, and that error is concentrated in
the smallest districts (the correlation between a district's standard
error and its log enrolment is -0.79).

\subsection{The predictors: socioeconomic status as six measures, not
one
index}\label{the-predictors-socioeconomic-status-as-six-measures-not-one-index}

Socioeconomic status could be operationalised as a single composite SEDA
ships one, \texttt{sesall}. We deliberately do not use it. The research
question asks \emph{which dimensions} of socioeconomic status carry the
most weight, and a composite answers that question by erasing it. So we
keep six separate measures: log median income, adult bachelor's-degree
share, poverty rate, unemployment, SNAP receipt, and single-mother
household share.

That choice has a cost, and it is worth quantifying rather than
glossing.

\begin{longtable}[]{@{}lr@{}}
\caption{Variance inflation factors, all six predictors entered
together.}\tabularnewline
\toprule\noalign{}
Variable & VIF \\
\midrule\noalign{}
\endfirsthead
\toprule\noalign{}
Variable & VIF \\
\midrule\noalign{}
\endhead
\bottomrule\noalign{}
\endlastfoot
Households below the poverty line (\%) & 6.15 \\
Log median household income & 6.11 \\
Households receiving SNAP (\%) & 4.81 \\
Adults with a bachelor's degree or higher (\%) & 2.98 \\
Households headed by a single mother (\%) & 2.72 \\
Unemployment rate (\%) & 1.89 \\
\end{longtable}

The six measures overlap heavily: poverty correlates with SNAP receipt
at 0.85 and with log income at -0.82. The variance inflation factors
show the consequence, and two predictors make it concrete. Fitted alone,
SNAP receipt carries a slope of -3.5 and log income 0.82, both
significant below 0.001. Entered alongside the other four, SNAP
\textbf{flips sign} to 0.2 (p = 0.64) while income collapses to 0.03 (p
= 0.8).

Neither variable has stopped mattering; each has been absorbed by its
correlates. We therefore \textbf{exclude SNAP receipt and log median
income} from the preferred specification not because they are
unimportant, but because their coefficients in a six-predictor model
would not mean what a reader would take them to mean. Poverty is
retained over SNAP because it is the more widely understood measure;
education and poverty together carry most of what income would have
contributed.

\subsection{Observations removed}\label{observations-removed}

The counts and reasons for every cleaning stage are given in the Data
Wrangling section above; the drop that matters is the \textbf{1,551
districts SEDA suppressed} for having too few tested students, plus a
further 95 missing one or more covariates. Because SEDA still publishes
covariates for the suppressed districts, we can check whether that loss
is systematic rather than assuming it. It is, but not in the direction
one would guess. Suppressed districts are genuinely smaller --- median
enrolment 391 against 561, and markedly so at the lower quartile --- and
somewhat more rural (59.6\% against 52.8\%). But they are \textbf{not}
poorer or less educated: mean poverty is 12.1\% against 12.7\%, and the
bachelor's-or-higher share 28.3\% against 24.3\%. The scope condition is
therefore narrower than a vague caveat would suggest: \textbf{our
results describe districts large enough for SEDA to publish a score},
and what is missing is very small, somewhat rural districts rather than
a systematically disadvantaged tail.

\section{The relationship between district poverty and test
performance}\label{the-relationship-between-district-poverty-and-test-performance}

This is the report's central exhibit: the outcome plotted against the
single predictor that carries most of the story, with the fitted model's
predictions overlaid.

\begin{figure}
\centering
\pandocbounded{\includegraphics[keepaspectratio,alt={District poverty and average test performance, 300 randomly sampled U.S. school districts, 2018. Each point is one district. The solid line is the fitted relationship from a model of test score on log poverty rate; the shaded band is its 95\% confidence interval. Districts are coloured by community type. The relationship is steep among low-poverty districts and flattens as poverty rises.}]{/Users/macandkash/Coding/dataci203-Final/datasci-203-26-section6-summer-C/notebooks/DATASCI-203-Final-Project_files/40fbf242a0d9343802faa4e95e25370e373131c9.pdf}}
\caption{District poverty and average test performance, 300 randomly
sampled U.S. school districts, 2018. Each point is one district. The
solid line is the fitted relationship from a model of test score on log
poverty rate; the shaded band is its 95\% confidence interval. Districts
are coloured by community type. The relationship is steep among
low-poverty districts and flattens as poverty rises.}
\end{figure}

the figure below displays all 300 sampled districts. The relationship is
strongly negative but plainly \textbf{not a straight line}: the fitted
curve falls steeply across the low-poverty range and flattens beyond
roughly 20\% poverty. Suburban districts cluster at the top left low
poverty, high scores while rural districts, the most numerous group at
163 of 300, spread across the full range. Community type is shown for
context only; it is not a term in any model.

Read as hypothetical districts, the curve says this. A district at 5\%
poverty is predicted to score +0.38 standard deviations; at 15\%, -0.14;
at 30\%, -0.47. The move from 5\% to 15\% is associated with a fall of
about 0.52 standard deviations, while the \emph{larger} move from 15\%
to 30\% is associated with only about 0.33. Ten points of poverty are
worth roughly twice as much at the bottom of the range as fifteen points
are at the top. That asymmetry is the descriptive finding, and it is why
the poverty rate enters the model in logs: a straight line would force
both moves to cost the same.

\section{Appendix: specifications
tried}\label{appendix-specifications-tried}

\begin{longtable}[]{@{}llrr@{}}
\caption{Three functional forms, tested on both variables where
curvature was plausible. Higher R-squared and lower residual SE are
better.}\tabularnewline
\toprule\noalign{}
Variable & Form & R-squared & Residual SE \\
\midrule\noalign{}
\endfirsthead
\toprule\noalign{}
Variable & Form & R-squared & Residual SE \\
\midrule\noalign{}
\endhead
\bottomrule\noalign{}
\endlastfoot
Poverty rate & Linear & 0.4481 & 0.2831 \\
Poverty rate & Quadratic & 0.4834 & 0.2744 \\
Poverty rate & Log & 0.4924 & 0.2715 \\
Bachelor's+ share & Linear & 0.4592 & 0.2803 \\
Bachelor's+ share & Quadratic & 0.4595 & 0.2807 \\
Bachelor's+ share & Log & 0.4372 & 0.2859 \\
\end{longtable}

{\def\LTcaptype{none} % do not increment counter
\begin{longtable}[]{@{}
  >{\raggedright\arraybackslash}p{(\linewidth - 2\tabcolsep) * \real{0.2743}}
  >{\raggedright\arraybackslash}p{(\linewidth - 2\tabcolsep) * \real{0.7257}}@{}}
\toprule\noalign{}
\begin{minipage}[b]{\linewidth}\raggedright
Specification
\end{minipage} & \begin{minipage}[b]{\linewidth}\raggedright
What we learned
\end{minipage} \\
\midrule\noalign{}
\endhead
\bottomrule\noalign{}
\endlastfoot
Y on poverty rate, untransformed & Fits, but the residual pattern and
the loess curve both show the straight line is wrong at the extremes. \\
Y on poverty rate plus a squared term & Curvature is real and strongly
significant, so the linear form is rejected. But a fitted parabola
eventually turns and predicts higher scores at very high poverty. \\
Y on log(poverty rate) & Best fit of the three, monotone across the
whole range, and interpretable as diminishing returns. Adopted. \\
Y on bachelor's share, log and squared variants & Neither improves on
the linear form, so education enters untransformed. Curvature is not a
general feature of these predictors. \\
Y on all six socioeconomic measures & SNAP receipt flips sign and log
income collapses. The measures are too collinear to interpret
individually. \\
Y on log(poverty), education, unemployment, single motherhood & Retains
the explanatory power of the six-predictor model with variance inflation
near or below three. Preferred specification. \\
\end{longtable}
}

The comparison table above is the evidence for rows one to four: for
poverty the log form fits best, while for the bachelor's-degree share
the linear form already does, so only poverty is transformed.

\begin{figure}
\centering
\pandocbounded{\includegraphics[keepaspectratio,alt={Residuals against fitted values for the proposed four-variable model. The local fit (orange) is close to flat with no funnelling, consistent with a correctly specified conditional mean and roughly constant error variance.}]{/Users/macandkash/Coding/dataci203-Final/datasci-203-26-section6-summer-C/notebooks/DATASCI-203-Final-Project_files/23cdb52cc3b312857a87738ba3b76051d3b268a7.pdf}}
\caption{Residuals against fitted values for the proposed four-variable
model. The local fit (orange) is close to flat with no funnelling,
consistent with a correctly specified conditional mean and roughly
constant error variance.}
\end{figure}

the residual plot below is the required residuals-versus-fitted plot for
the preferred specification. The local fit is close to flat with no
systematic curvature, which is consistent with a correctly specified
conditional mean had the poverty rate been left on its raw scale, the
curvature documented above would appear here as a parabolic band.

\section{Model Building}\label{model-building}

All model choices below are based only on the exploration data. We do
not use the confirmation set until the very end, so none of the
decisions will introduce bias in the final results.

We fit three models. \textbf{Model 1} regresses test performance on the
poverty rate alone, the simplest credible specification and the measure
policy discussions lean on most. \textbf{Model 2} adds the other five
socioeconomic factors, showing how much of poverty's apparent
association survives conditioning. \textbf{Model 3} adds a squared
poverty term to test whether the relationship bends at high poverty. All
three report heteroskedasticity-consistent standard errors, since
district-level data of this kind commonly has non-constant error
variance.

\begin{table}[!htbp] \centering 
  \caption{Regression Results (Exploration Set)} 
  \label{} 
\small 
\begin{tabular}{@{\extracolsep{5pt}}lccc} 
\\[-1.8ex]\hline 
\hline \\[-1.8ex] 
 & \multicolumn{3}{c}{\textit{Dependent variable:}} \\ 
\cline{2-4} 
\\[-1.8ex] & \multicolumn{3}{c}{Test Performance (SD units from national mean)} \\ 
\\[-1.8ex] & (1) & (2) & (3)\\ 
\hline \\[-1.8ex] 
 Poverty Rate & $-$4.02$^{***}$ & $-$1.24$^{*}$ & $-$2.59$^{*}$ \\ 
  & (0.29) & (0.57) & (1.18) \\ 
  Poverty Rate Squared &  &  & 4.14 \\ 
  &  &  & (3.09) \\ 
  Log Median Income &  & 0.03 & $-$0.01 \\ 
  &  & (0.12) & (0.12) \\ 
  \% Bachelor's Degree+ &  & 1.30$^{***}$ & 1.26$^{***}$ \\ 
  &  & (0.18) & (0.19) \\ 
  Unemployment Rate &  & $-$2.30$^{*}$ & $-$2.38$^{*}$ \\ 
  &  & (1.10) & (1.11) \\ 
  \% SNAP Recipients &  & 0.20 & 0.14 \\ 
  &  & (0.48) & (0.48) \\ 
  \% Single-Mother Households &  & $-$1.19$^{***}$ & $-$1.23$^{***}$ \\ 
  &  & (0.36) & (0.35) \\ 
  Constant & 0.51$^{***}$ & $-$0.16 & 0.32 \\ 
  & (0.04) & (1.27) & (1.32) \\ 
 \hline \\[-1.8ex] 
Observations & 300 & 300 & 300 \\ 
R$^{2}$ & 0.45 & 0.63 & 0.63 \\ 
Adjusted R$^{2}$ & 0.45 & 0.62 & 0.62 \\ 
\hline 
\hline \\[-1.8ex] 
\textit{Note:}  & \multicolumn{3}{r}{$^{*}$p$<$0.05; $^{**}$p$<$0.01; $^{***}$p$<$0.001} \\ 
\end{tabular} 
\end{table}

\subsection{Assumption checks}\label{assumption-checks}

With 300 districts the large-sample model needs only IID sampling and a
unique BLP, both of which the sampling design gives us directly. The
supporting diagnostics agree: a studentized Breusch--Pagan test returns
p = 0.285, giving no evidence of heteroskedasticity, and while
Shapiro--Wilk rejects exact normality of the residuals (p = 0.0043) the
excess kurtosis is only 1.23, which at this sample size does not
threaten the asymptotic results. We report robust standard errors
regardless. The residuals-versus-fitted plot for the preferred
specification is Figure 2 in the appendix.

\end{document}
